\section{Method}
\label{sec:method}

We propose DFM-Fusion (Deterministic FadeMem-Fusion), a fully deterministic replacement for LLM-guided memory fusion in long-horizon conversational agents. Our approach maintains the memory architecture and retrieval mechanisms of existing systems while eliminating all LLM calls from the fusion operator.

\subsection{Problem Formulation}

Memory fusion addresses the challenge of consolidating redundant information in conversational memory stores. Given a cluster of semantically similar memory entries $\mathcal{M} = \{m_1, m_2, \ldots, m_k\}$ identified through temporal-semantic clustering, the fusion task produces a single consolidated entry $m_f$ that preserves essential information within a token budget $B$.

Existing approaches~\citep{Wei2026FadeMemBF,chhikara2025mem0buildingproductionreadyai} employ LLMs to rewrite and merge memories, then use an additional LLM call to verify that the rewrite preserved information. As discussed in Section~\ref{sec:intro}, this introduces deployment challenges related to cost, determinism, and auditability. Our approach addresses these by preserving verbatim spans while enforcing explicit token budgets.

\paragraph{Memory strength.}
We follow the FadeMem-style strength dynamics used by the implementation. Each new memory starts with strength $v_i=1.0$ in the short-term layer. At time $t$, its importance is
\begin{equation}
I_i(t)=\alpha r_i(t)+\beta \frac{a_i}{1+a_i}+\gamma \exp(-\delta \Delta t_i),
\end{equation}
where $r_i(t)$ is cosine relevance to the current session context when available, $a_i$ is access frequency, and $\Delta t_i$ is age in days. Strength decays as
\begin{equation}
v_i(t)=v_i(t_0)\exp\left[-\lambda_\text{base}\exp(-\mu I_i(t))(\Delta t_i)^{\rho}\right],
\end{equation}
with $\rho=0.8$ for long-term memories and $\rho=1.2$ for short-term memories. Access consolidation increases strength by
\begin{equation}
v_i \leftarrow v_i + d_v(1-v_i)\exp(-c_i/N),
\end{equation}
where $c_i$ is access count. A fused memory receives strength
\begin{equation}
v_f=\min\left(1,\max_i v_i+\epsilon \operatorname{Var}(\{v_i\})\right),
\end{equation}
and is inserted into the long-term layer. In our experiments $\alpha=0.4$, $\beta=0.3$, $\gamma=0.3$, $\delta=0.1$, $\mu=2.0$, $d_v=0.15$, $N=5$, and $\epsilon=0.1$.

\subsection{DFM-Fusion Pipeline}

Given a cluster $\mathcal{M}$ exceeding a minimum size threshold, DFM-Fusion produces fused text $s_f$ through four deterministic stages. More specifically: (1) sentence segmentation breaks each memory entry into sentence-level units to enable fine-grained selection; (2) near-duplicate removal eliminates redundant sentences via embedding similarity; (3) MMR-based sentence packing selects a diverse, high-strength subset within the token budget; and (4) coverage verification confirms that salient evidence is preserved, with a fallback guarantee if it is not. No LLM is invoked at any stage.

\paragraph{Stage 1: Sentence Segmentation.}
Each memory entry $m_i$ is split into individual sentences using rule-based tokenization. However, applying fusion directly at the memory-entry level would truncate content mid-sentence when the token budget is enforced, fragmenting the verbatim spans that downstream retrieval and F1 evaluation depend on. Sentence-level segmentation ensures that every selected unit is a complete, citable span.

\paragraph{Stage 2: Near-Duplicate Removal.}
We compute pairwise cosine similarity between sentence embeddings using Sentence-BERT~\citep{Reimers2019SentenceBERTSE} (384-dim). However, the raw sentence pool from a cluster frequently contains near-identical reformulations of the same fact --- variants that would consume token budget without adding evidence coverage. Sentences with cosine similarity exceeding threshold $\theta_{\text{dup}} = 0.85$ are treated as near-duplicates; we retain the sentence from the memory with higher strength $v_i(t)$, using a more recent timestamp as a tie-breaker. The same Sentence-BERT encoder is reused throughout the pipeline --- for deduplication, for MMR scoring, and for computing the retrieval embedding of the final fused entry --- eliminating any inconsistency between maintenance-time and retrieval-time representations.

\paragraph{Stage 3: MMR-Based Sentence Packing.}
From the deduplicated sentence pool, we select sentences using a strength-aware variant of Maximal Marginal Relevance~\citep{Carbonell1998MMR} with diversity parameter $\lambda = 0.7$. However, relevance scoring alone would fill the budget with the highest-strength sentences from the dominant source memory, crowding out lower-strength entries that may contain complementary evidence. The MMR objective balances source-memory strength against redundancy with already-selected sentences:
\begin{equation}
\text{MMR}(s_i) = \lambda \cdot \hat{v}(s_i) - (1-\lambda) \cdot \max_{s_j \in S} \cos(s_i, s_j)
\end{equation}
where $\hat{v}(s_i) \in [0,1]$ is the normalized strength of the source memory for sentence $s_i$, and $S$ is the set of already-selected sentences. We greedily pick $\arg\max_{s_i \in R \setminus S} \text{MMR}(s_i)$ until adding the next sentence would exceed the token budget $B_{\text{fuse}} = 768$ tokens. Selected sentences are sorted by source-memory timestamp before concatenation. The fused memory embedding is computed from the actual fused text $s_f$, not by averaging constituent embeddings; this ensures retrieval operates on a representation that accurately reflects the consolidated content, and was validated by our ablation (Table~\ref{tab:main_results}, lower block).

\paragraph{Stage 4: Coverage Verification.}
To confirm information preservation without an LLM, we implement a deterministic coverage check. We extract salient evidence from the original cluster: numbers and date-like expressions, capitalized entities, top-$K$ TF-IDF weighted tokens ($K=20$), and stemmed $n$-gram coverage features in the enriched configuration. Coverage recall is the fraction of these tokens appearing in $s_f$. It should be noted that a coverage gate that simply rejects low-recall fusions would leave clusters unmerged, breaking the consolidation guarantee. To address this, when recall falls below threshold $\theta_{\text{cov}} = 0.85$, we fall back to concatenating the highest-strength original memories within the budget in strength-descending order, so a fused result is always emitted rather than leaving the cluster unmerged.

\subsection{DFM-Route: Query-Aware Retrieval Extension}
\label{sec:dfm-route}

While DFM-Fusion operates at maintenance time --- consolidating memory clusters offline before any query arrives --- DFM-Route extends the system at retrieval time by routing each query to a retrieval strategy matched to its evidence requirements.
Like DFM-Fusion, DFM-Route introduces zero LLM calls at any stage; all routing and verification decisions are fully deterministic.

\paragraph{Query Type Classification.}
Each incoming query is classified into one of eight types by a rule-based regex classifier: \textit{temporal} (when-questions), \textit{list}, \textit{number}, \textit{who}, \textit{where}, \textit{yes\_no}, \textit{inferential\_yes\_no}, or \textit{fact} (default).
Classification is based solely on surface patterns (e.g., leading ``when'', ``who'', ``how many'') without any learned model.

\paragraph{Evidence Gate Score.}
For each query, two candidate retrieval results are generated: a \textit{conservative} set (standard cosine-similarity ranking) and an \textit{expansive} set (anchor-then-MMR, which locks on the highest-similarity memory and applies MMR diversification).
A deterministic gate score $g \in [0,1]$ is computed for each candidate set:
\begin{equation}
    g = 0.35\,\sigma_{\max} + 0.30\,\rho_{\text{unit}} + 0.20\,\rho_{\text{subj}} + 0.10\,a + 0.05\,b,
    \label{eq:gate}
\end{equation}
where $\sigma_{\max}$ is the maximum embedding similarity in the retrieved set, $\rho_{\text{unit}}$ is the overlap between query key units (entities, dates, numerics) and retrieved tokens, $\rho_{\text{subj}}$ is subject-phrase overlap, $a$ is the Jaccard agreement between the conservative and expansive sets, and $b = \min(1, |\text{sources}|/6)$ is a source-diversity bonus.

\paragraph{Routing Decision.}
The route is selected by comparing the gate scores of the two candidate sets against type-specific thresholds.
For \textit{temporal} queries, the conservative set is used by default (\textit{temporal guard}) unless the expansive gate score exceeds $\theta_{\text{temp}}=0.62$ with agreement $\geq 0.25$, in which case the expansive set is used.
For \textit{list}, \textit{number}, \textit{who}, and \textit{where} queries, the \textit{union} of both sets is used when subject overlap is sufficient and gate scores meet minimum thresholds.
For \textit{fact} queries, the expansive set is used only when it shows a clear gate-score advantage ($\Delta g \geq 0.08$) over the conservative set; otherwise the conservative set is used.
Across 1,986 questions on LoCoMo, this yields 78\% conservative, 12\% union, and 10\% expansive routing.

\paragraph{Rule-Based Answer Verifier.}
After answer generation, a lightweight verifier checks whether the answer is supported by the retrieved context.
The verifier applies type-specific rules without any LLM: for yes/no questions it checks substring support; for list questions it verifies each item against the retrieved evidence; for temporal questions it attempts date-unit recovery from the context.
In practice the verifier keeps 97\% of answers unchanged, recovers 2\%, and abstains on 1\%, confirming that its role is conservative post-hoc filtering rather than aggressive rewriting.

\subsection{Theoretical Analysis}
\label{sec:theory}

We establish two formal properties that together justify DFM-Fusion's design and explain its empirical advantage on fact-centric benchmarks.
The first (Lemma~\ref{lem:submodular}) shows that coverage maximization is a submodular optimization problem, justifying the greedy sentence selection in Stage~3.
The second (Theorem~\ref{thm:equiv}) identifies the structural condition under which verbatim packing is provably equivalent to LLM-guided rewriting in terms of answer recall.

\paragraph{Coverage maximization as submodular optimization.}
Let $\mathcal{S} = \bigcup_{i=1}^{k} \mathrm{sent}(m_i)$ be the full sentence pool after Stage~1, and let $T = T_{\mathrm{num}} \cup T_{\mathrm{ent}} \cup T_{\mathrm{tfidf}} \cup T_{\mathrm{ng}}$ be the salient token set assembled in Stage~4.
Define the coverage function
\begin{equation}
    f(A) = \frac{\bigl|T \cap \textstyle\bigcup_{s \in A} \mathrm{tok}(s)\bigr|}{|T|}, \quad A \subseteq \mathcal{S},
    \label{eq:coverage}
\end{equation}
where $\mathrm{tok}(s)$ denotes the token multiset of sentence $s$.
The DFM-Fusion selection problem is
\begin{equation}
    \max_{A \subseteq \mathcal{S}} \; f(A) \quad \text{s.t.} \quad \textstyle\sum_{s \in A}|s|_{\mathrm{tok}} \leq B_{\mathrm{fuse}}.
    \label{eq:opt_fusion}
\end{equation}

\begin{lemma}[Submodularity of Coverage]
\label{lem:submodular}
The function $f$ defined in~\eqref{eq:coverage} is monotone submodular.
\end{lemma}
\begin{proof}
\emph{Monotonicity}: For $A \subseteq A'$, $\bigcup_{s \in A}\mathrm{tok}(s) \subseteq \bigcup_{s \in A'}\mathrm{tok}(s)$, so $f(A) \leq f(A')$.
\emph{Diminishing returns}: For any $A \subseteq B \subseteq \mathcal{S}$ and $s \notin B$, the marginal gain satisfies
\[
  f(B \cup \{s\}) - f(B)
  = \frac{\bigl|T \cap \mathrm{tok}(s) \setminus \textstyle\bigcup_{t \in B}\mathrm{tok}(t)\bigr|}{|T|}
  \;\leq\;
  \frac{\bigl|T \cap \mathrm{tok}(s) \setminus \textstyle\bigcup_{t \in A}\mathrm{tok}(t)\bigr|}{|T|}
  = f(A \cup \{s\}) - f(A),
\]
since $A \subseteq B$ implies $\bigcup_{t \in A}\mathrm{tok}(t) \subseteq \bigcup_{t \in B}\mathrm{tok}(t)$, so the tokens excluded from $s$'s marginal contribution are a superset under $B$.
\end{proof}

Because $f$ is monotone submodular and the budget constraint is a modular knapsack, a greedy algorithm that iteratively selects the sentence with the highest marginal coverage-per-token achieves a $(1-1/e)$-approximation to the optimal solution of~\eqref{eq:opt_fusion}~\citep{Nemhauser1978AnalysisOA,Sviridenko2004ANO}.
Stage~3 maximizes a strength-weighted generalization of $f$ via greedy MMR; the submodular structure is preserved because the strength weights enter as non-negative scalars on the marginal coverage terms.

\paragraph{When verbatim packing is sufficient.}

\begin{definition}[Verbatim Sufficiency]
\label{def:vs}
A QA instance $(q, a^*)$ is \emph{Verbatim Sufficient} (VS) with respect to cluster $\mathcal{M}$ if every token of the gold answer $a^*$ appears verbatim in at least one sentence of the pool:
\[
  \forall\, w \in \mathrm{tok}(a^*),\quad \exists\, s \in \mathcal{S} : w \in \mathrm{tok}(s).
\]
\end{definition}

\begin{theorem}[Recall Dominance under VS]
\label{thm:equiv}
Let $(q, a^*)$ be VS with respect to $\mathcal{M}$, and let $s_f = \mathrm{concat}(A)$ be the DFM-Fusion output such that $\mathrm{tok}(a^*) \subseteq \mathrm{tok}(s_f)$ (the answer-bearing sentences are included within budget).
Then for any rewriting-based fusion output $s_r$ of $\mathcal{M}$:
\[
  \mathrm{Recall}(s_f,\, a^*) = 1 \;\geq\; \mathrm{Recall}(s_r,\, a^*),
\]
where $\mathrm{Recall}(s, a^*) = |\mathrm{tok}(s) \cap \mathrm{tok}(a^*)| \,/\, |\mathrm{tok}(a^*)|$.
\end{theorem}
\begin{proof}
By assumption, $\mathrm{tok}(a^*) \subseteq \mathrm{tok}(s_f)$, so $|\mathrm{tok}(s_f) \cap \mathrm{tok}(a^*)| = |\mathrm{tok}(a^*)|$, giving $\mathrm{Recall}(s_f, a^*) = 1$.
Since $\mathrm{Recall} \in [0,1]$, the inequality $1 \geq \mathrm{Recall}(s_r, a^*)$ is immediate for any $s_r$.
\end{proof}

Theorem~\ref{thm:equiv} establishes that under VS, DFM-Fusion achieves perfect answer-token recall---the maximum possible---making LLM-guided rewriting redundant for recall on fact-centric questions.
Any LLM rewrite that paraphrases or omits answer tokens strictly reduces recall below 1; any rewrite that preserves them achieves recall equal to DFM-Fusion.
The condition ``answer-bearing sentences within budget'' is satisfied when the budget $B_{\mathrm{fuse}}$ is large enough to accommodate the salient high-strength members that carry the answer, as our Stage~3 MMR selection prioritizes.

\begin{remark}
\label{rem:vs_boundary}
VS holds for fact-centric questions whose answers are directly stated in the conversation: entity names, precise dates, and numerical facts.
It fails when the gold answer requires synthesis not present verbatim in the cluster --- for example, normalizing ``she joined three months after moving'' into a relative temporal expression.
This boundary is consistent with DFM-Fusion's observed advantages on multi-hop and single-hop categories and LLM-Fusion's retained advantages on temporal and open-domain categories (Table~\ref{tab:main_results}).
\end{remark}

\begin{remark}[Coverage Certificate vs.\ LLM Preservation Check]
\label{rem:cert}
Stage~4 provides a \emph{certified} lower bound: $f(A) \geq \theta_{\mathrm{cov}}$ guarantees that at least $\theta_{\mathrm{cov}}$ fraction of salient tokens appear verbatim in $s_f$.
By contrast, LLM-Fusion's preservation check produces an uncertified judgment that may hallucinate confirmation or miss subtle omissions.
The certificate is weaker in scope (token-level coverage, not semantic equivalence) but strictly stronger in reliability: it is deterministic and auditable.
\end{remark}

\subsection{Hyperparameter Selection}

FadeMem decay parameters ($\theta_\text{fusion}=0.75$, $T_\text{window}=10$, $\lambda_\text{base}=0.1$) are adopted directly from the original system. DFM-specific fusion parameters were initialized to interpretable defaults ($\theta_\text{dup}=0.90$, $B_\text{fuse}=512$, $\lambda=0.5$, $K=10$, $\theta_\text{cov}=0.80$) and revised once, in a single diagnostic pass that identified four concrete issues: an embedding representation bug (weighted-average rather than fused-text encoding), a suboptimal text separator, an over-aggressive deduplication threshold, and an undersized fusion budget. This pass was conducted on the full LoCoMo-10 evaluation set without a held-out split. While the changes are individually motivated by interpretable diagnostics rather than metric search, we acknowledge that we cannot fully rule out mild overfitting to the 10 evaluation conversations; the six-model cross-family replication (Section~\ref{sec:experiments}) provides indirect stability evidence.
